\documentclass[11pt,letterpaper]{article}
\usepackage[margin=1in]{geometry}
\usepackage[T1]{fontenc}
\usepackage{newtxtext,newtxmath,setspace,xurl}
\usepackage[hidelinks]{hyperref}
\onehalfspacing\setlength{\emergencystretch}{3em}
\title{Response to the Independent R45 Referee Report}
\author{Anonymous manuscript for Operations Research}
\date{September 25, 2026}
\begin{document}\maketitle
\section*{Version and scope}
This response addresses the independent R45 report dated September 25, 2026, on review commit 260a5224c55e0326d9da7f0c813b4a3fcab7671f. That report reviewed R45 manuscript tip 9c52e2d58d69a8d653c3b50e2c860b35db04dc7f. R46 is based on the review commit and is confined to revision/ndu-operations-research-r46-box-decomposition-20260925. The original manuscript and every historical derivation are retained; the revision does not change the original alphabet budget, opening-charge convention, root equality, or realized ceilings.

\section*{Main change: a box oracle rather than a renamed target grid}
We agree that completing the prior target grid is not a scalable joint-design mechanism in arbitrary branch dimension, and that fixed-target timing cannot establish otherwise. We have replaced that grid as the main computational argument with an exact two-sided price oracle on arbitrary target boxes, an adaptive joint algorithm with local tightness, and a separately checked cover certificate. The prior grid remains with its original theorem and proof in the companion. The method is not claimed to dominate every inherited algorithm. The reported failures and the stronger performance of price-prefix search on several instances remain in the revised record.

\section*{R1. Complexity frontier and repeated branch types}
Propositions on exact response-type reduction and positive complexity regimes give a precise parameter map. Histories with the same cap, cost function, and eligible catalog prefix aggregate exactly; different numerical ceilings are permitted when their eligible prefixes agree. This extends, rather than relabels, the earlier homogeneous common-ceiling aggregation argument. One response type and saturation each reduce to one conditional path optimization. For every fixed alphabet budget, exact book enumeration followed by concave allocation is polynomial in catalog size and input bit length. This is an XP statement in the alphabet budget, not an FPT claim. The adaptive accuracy bound uses the number of distinct response types minus one, rather than raw history count. We do not claim a general NP-hardness classification or a dimension-independent polynomial algorithm. In particular, we do not turn the old grid complexity into an unsupported hardness lower bound.

\section*{R2. Structural global method}
The two-sided box-price theorem is the main new structural result. A positive-price ordered book has an increasing prefix and a decreasing suffix. Branches with upper targets below the peak contribute upper-crossing edges; branches whose target intervals contain the peak contribute its score; branches forced above the peak contribute lower-crossing edges. Coupled forward and backward recurrences optimize this decomposition in polynomial work per node and price, preserving every selected-level charge and the original cardinality. This is not obtained by calling the old unrestricted price oracle on a smaller box. We supply an exact three-target counterexample: the rising-only restriction returns 83/100 while the full box value is 169/200, a deficit of 3/200. The new recurrence attains the correct value.

\section*{R3. Global accuracy and independent coverage}
The adaptive theorem proves a valid original-budget interval at every stopping point and finite termination at every positive tolerance. A root-feasible anchor condition ensures that the zero-price maximizing book also admits an equality-feasible allocation in the box. Lipschitz response bounds then make the node upper bound locally tight. The proof supplies a conservative dyadic depth bound and explicitly retains exponential worst-case dependence on the number of types. A separate implementation checks every forward and backward Bellman inequality, all projected child boxes, every authoritative leaf, and the original uneliminated lottery. It does not import the search routine or accept coverage from a completion flag. The prior grid is also tested at tolerances 0.1, 0.01, and 0.001 with two through five branches; incomplete grids explicitly withhold their requested-accuracy certificate.

\section*{R4. Joint experiments versus heuristics}
The new primary result concerns 32 joint-design cases, not a conditional frontier or the R45 target portfolio. Twenty-eight reach absolute gap 0.001; four retain nonzero certified gaps and no requested-accuracy guarantee. All cases, including interruption, are in the tables and machine-readable records. The cases use fixed independent seeds; protocols were frozen locally before their respective executions, not externally preregistered. The 128- and 256-history cases explicitly use repeated types and check the lifted policy on the original history space. These examples demonstrate exact type reduction, not evidence that arbitrary operational populations have four types. The older 35/36 and safeguarded 36/36 portfolio counts remain historical observations, with corrected headings, rather than an algorithmic headline.

\section*{R5. Independent comparisons and matched objectives}
We compare the new method with the inherited R43 priced-prefix completion bound, reimplemented at the same requested accuracy and root price budget; a simpler full-catalog prefix bound is additionally retained and is not substituted for the stronger baseline. Price-prefix search reaches the requested tolerance on all 32 cases, including the four on which target-box refinement stops early. The new box method is faster on 11 measured cases, but there is no claim of universal superiority. Every method keeps the same finite catalog, cardinality, charges and promise. The comparison has an eight-second nominal allowance, with documented soft checks and separate node guards. The direct uneliminated MIP is run on all 24 primary cases with tangent and secant cost envelopes, 64 segments and a four-second allowance per solve. All 48 solves returned optimal numerical status here; per-case times, nodes, mixed-integer gaps, approximation error, bracket widths and residuals are retained. Numerical bounds remain numerical even when they match rational references.

\section*{R6. Stronger large-instance certificates and inexact supports}
Every box bound is target-sensitive and optimizes the shared alphabet, rather than evaluating Jensen at one supplied target vector. The numerical section reports the improvement from the Jensen-minus-incumbent interval to the final box interval and retains instances where the latter remains too wide. The inexact-support composition theorem is preserved in the companion with its assumptions. The new global interval does not require an externally supplied certified global inexact price oracle: for the rational quadratic class the box oracle computes its own exact upper bound. It also does not require mixing two endpoint books or bounding a compression penalty. General cost functions still require the stated scalar support oracle; arbitrary approximate cost oracles are not silently treated as exact.

\section*{R7. Application scope, presentation, and preservation}
The revised main argument is an optimization contribution: a lower-bound-aware price decomposition and a certified adaptive joint method. We do not fabricate a calibrated operational case. The main text now distinguishes terminal command bits from the history encoder, target and lottery tables, read-only constants, probability representation, and offline certificate memory. Canonical institutional, continuous, Monge, catalog, augmentation, and recovery results retain their mathematical statements and proofs across the current article and companion. Old numerical protocols and full predecessor readers are preserved in the evidence archive with an explicit content map, rather than being presented again as new executions. The article uses anonymous authorship, 11-point text, one-inch margins, one-and-a-half spacing, an abstract below 200 words, a nonmathematical introduction, author-year citations, and tables after references. The build records page counts and checks both the lengthy-manuscript and companion-size limits.

\section*{M1. One-branch target-net redundancy}
The new target\_net\_reference.py returns the exact conditional frontier after one call at the forced target. The regression explicitly checks evaluated\_profiles equals one. The archived R45 implementation is not overwritten.

\section*{M2. Clipping regression}
The repair tests now include candidate coordinates below the lower anchor and above the cap, and both upward and downward aggregate repair. They assert strictly positive clipping cost as well as exact final equality and box feasibility.

\section*{M3. Coverage checker}
The new independent checker validates the complete adaptive cover, not merely sampled target agreement. We additionally execute 2,793 finite-lattice constructive coverage checks for the retained net. Those finite regressions support implementation verification and do not replace its universal proof.

\section*{M4. Grid timing, scale, and loss}
The new grid table identifies catalog size and budget, runtime, profile count, tolerance and completion. grid.json additionally records the exact objective, absolute loss, positive declared normalization scale and both algorithms' times. A profile-limited run is not described as a time-matched speedup experiment.

\section*{M5. Saturated algorithm provenance}
The theorem-level comparison table identifies the charged saturated path as an R42 antecedent generalized by the R45 fixed-target recurrence. Neither is counted as the new R46 box theorem.

\section*{M6. Title and joint scope}
The title is now Limited-Memory Renewal Contracts: Two-Sided Decomposition and Certified Joint Design. Its joint-design claim is supported by the box interval and finite-accuracy theorem, not by large fixed-target timings. Performance and worst-case qualifications remain explicit.

\section*{M7. Exponential accuracy dependence}
The abstract directly states that worst-case accuracy dependence is exponential in the number of distinct types. The theorem gives the dyadic bound and does not call it an FPTAS for variable dimension.

\section*{M8. Safeguarded sample label}
The reader copy of the historical table now says Safeguarded exact in sample. Every historical row and its unaltered source remain preserved.

\section*{M9. Charged counterexample scale}
For the retained seed 205019, budget two, the core portfolio value is 53/150 and the exact optimum is 1373/3600. Their difference is 101/3600, about 7.36 percent of that positive exact net optimum. The failure is not omitted or relabeled as roundoff.

\section*{M10. MIP reporting}
The companion table covers every one of the 24 new primary comparisons. The raw records give both envelope statuses, runtimes, node counts, primal--dual gaps, envelope errors, solutions, residuals and signed bracket widths. All solves completed here; no failed or timed-out case was filtered. Tiny negative signed widths at floating roundoff scale are preserved, not called exact intervals.

\section*{M11. Inserted ideal targets}
The historical interior target-adaptive row is relabeled oracle pool, and the caption states that ideal targets were inserted. It is a sanity check of an attainable bound, not evidence for a general adaptive candidate algorithm.

\section*{M12. Conditional scaling distinction}
The historical caption explicitly says targets fixed. The new joint table instead reports an upper and lower interval for target and book optimization. The two sets are not pooled.

\section*{M13. Finite-catalog scope}
The abstract and new statements explicitly identify the finite-catalog problem. The separate continuous quadratic and mesh-transfer results retain their own assumptions; no unrestricted continuous guarantee is inferred from a finite catalog.

\section*{M14. Theorem-level comparison}
The main comparison table distinguishes canonical and continuous antecedents, the R42 saturated algorithm, R43 price decomposition and search, R45 conditional compression and target net, and the R46 type, box, adaptive and checker contributions.

\section*{M15. Selected versus used levels}
The text defines the budget as selected enabled levels, not positive-probability support size. Historical Used is relabeled Selected. Full original-space lotteries retain enough information to compute actual positive-probability usage separately.

\section*{M16. Additive scale}
The experiments state normalized target and reward primitives and report gaps relative to a positive declared objective scale even when net value is near zero. Rescaling reward and opening charges requires rescaling the absolute tolerance.

\section*{M17. Repeated types}
The exact type proposition and implemented grouping precede the joint search. The equality of aggregated and disaggregated exact optima is regression-tested, including different numerical ceilings with the same catalog prefix. The implementation does not merge branches with merely similar caps.

\section*{M18. Numerical versus rational evidence}
All target-box intervals and their checks use exact rational arithmetic. Floating-point HiGHS output remains separately labeled. Source hashes, solver status, sample exactness and a feasible policy are not substituted for a global mathematical certificate.

\section*{Remaining qualifications}
The revision supplies new structural theory and an implemented globally certifying algorithm, rather than a plan to add them later. It does not claim a complete general hardness classification, uniformly superior runtime, or a field-calibrated application. The four unresolved primary box runs and the interrupted four-branch fine-grid comparison remain visible. These limits qualify computational conclusions without changing the original optimization problem or removing the contractual and continuous contributions.

\end{document}
